\documentclass[12pt]{article}
\usepackage[a4paper,margin=1in]{geometry}
\usepackage[utf8]{inputenc}
\usepackage[T1]{fontenc}
\usepackage{amsmath,amssymb}
\usepackage{textcomp}
\usepackage{parskip}
\begin{document}
\begin{center}
{\Large \textbf{1998 STEP I}}
\end{center}

\subsection*{Page 1}
\begin{flushleft}
\bigskip\textbf{Section A: Pure Mathematics}\par
\bigskip\textbf{Question 1.} How many integers between 10 000 and 100 000 (inclusive) contain exactly two different digits?\par
(23 332 contains exactly two different digits but neither of 33 333 and 12 331 does.)\par
2 Show, by means of a suitable change of variable, or otherwise, that\par
\infty \infty\par
f((x\par
+ 1)\par
\[1/2\]
\[+ x) dx = (1 + t\]
-2\par
)f(t) dt.\par
\[Z^0 2 Z^1\]
Hence, or otherwise, show that\par
\infty\par
((x\par
+ 1)\par
\[1/2\]
+ x)\par
-3\par
\[dx = .\]
\[Z^0 8\]
\bigskip\textbf{Question 3.} Which of the following statements are true and which are false? Justify your answers.\par
\textbf{(i)} aln b = bln a for all a, b > 0.\par
\textbf{(ii)} cos(sin θ) = sin(cos θ) for all real θ.\par
\textbf{(iii)} There exists a polynomial P such that P(θ) cos θ 6 10 -6 for all real θ.\par
| - |\par
\textbf{(iv)} x4 + 3 + x -4 > 5 for all x > 0.\par
\bigskip\textbf{Question 4.} Prove that the rectangle of greatest perimeter which can be inscribed in a given circle is a\par
square.\par
The result changes if, instead of maximising the sum of lengths of sides of the rectangle, we\par
seek to maximise the sum of nth powers of the lengths of those sides for n > 2. What happens\par
\[if n = 2? What happens if n = 3? Justify your answers.\]
\end{flushleft}
\subsection*{Page 2}
\begin{flushleft}
5 (i) In the Argand diagram, the points Q and A represent the complex numbers 4 + 6i and\par
10 + 2i. If A, B, C, D, E, F are the vertices, taken in clockwise order, of a regular\par
hexagon (regular six-sided polygon) with centre Q, find the complex number which\par
represents B.\par
\textbf{(ii)} Let a, b and c be real numbers. Find a condition of the form Aa + Bb + Cc = 0, where\par
A, B and C are integers, which ensures that\par
a b c\par
+ +\par
1 + i 1 + 2i 1 + 3i\par
is real.\par
\bigskip\textbf{Question 6.} Let a1 = cos x with 0 < x < \pi/2 and let b1 = 1. Given that\par
\[a n+1 = 2 (a n + b n ),\]
\[b n+1 = (a n+1b n )\]
\[1/2\]
,\par
\[find a^2 and b^2 and show that\]
x x x x\par
\[a^3 = cos cos and b^3 = cos cos .\]
2 4 2 4\par
Guess general expressions for a and b (for n > 2) as products of cosines and verify that they\par
n n\par
satisfy the given equations.\par
\bigskip\textbf{Question 7.} My bank pays ρ\% interest at the end of each year. I start with nothing in my account. Then\par
for m years I deposit \pounds a in my account at the beginning of each year. After the end of the mth\par
year, I neither deposit nor withdraw for l years. Show that the total amount in my account\par
at the end of this period is\par
rl+1(rm 1)\par
\pounds a -\par
r 1\par
-\par
ρ\par
\[where r = 1 + .\]
100\par
At the beginning of each of the n years following this period I withdraw \pounds b and this leaves\par
my account empty after the nth withdrawal. Find an expression for a/b in terms of r, l, m\par
and n.\par
\end{flushleft}
\subsection*{Page 3}
\begin{flushleft}
\bigskip\textbf{Question 8.} Fluid flows steadily under a constant pressure gradient along a straight tube of circular cross-\par
section of radius a. The velocity v of a particle of the fluid is parallel to the axis of the tube\par
and depends only on the distance r from the axis. The equation satisfied by v is\par
1 d dv\par
\[r = k,\]
r dr µ dr ¶ -\par
where k is constant. Find the general solution for v.\par
Show that v as r 0 unless one of the constants in your solution is chosen to be 0.\par
| | \to \infty \to\par
Suppose that this constant is, in fact, 0 and that v = 0 when r = a. Find v in terms of k, a\par
and r.\par
The volume F flowing through the tube per unit time is given by\par
a\par
\[F = 2\pi rv dr.\]
\[Z^0\]
Find F .\par
\end{flushleft}
\subsection*{Page 4}
\begin{flushleft}
\bigskip\textbf{Section B: Mechanics}\par
\bigskip\textbf{Question 9.} Two small spheres A and B of equal mass m are suspended in contact by two light inextensible\par
strings of equal length so that the strings are vertical and the line of centres is horizontal. The\par
coefficient of restitution between the spheres is e. The sphere A is drawn aside through a very\par
small distance in the plane of the strings and allowed to fall back and collide with the other\par
sphere B, its speed on impact being u. Explain briefly why the succeeding collisions will all\par
occur at the lowest point. (Hint: Consider the periods of the two pendulums involved.)\par
Show that the speed of sphere A immediately after the second impact is 1 u(1 + e2) and find\par
the speed, then, of sphere B.\par
\bigskip\textbf{Question 10.} A shell explodes on the surface of horizontal ground. Earth is scattered in all directions with\par
varying velocities. Show that particles of earth with initial speed v landing a distance r from\par
the centre of explosion will do so at times t given by\par
1 2 2 2 4 2 2\par
\[g t = v \sqrt(v g r ).\]
2 ± -\par
Find an expression in terms of v, r and g for the greatest height reached by such particles.\par
\bigskip\textbf{Question 11.} Hank's Gold Mine has a very long vertical shaft of height l. A light chain of length l passes\par
over a small smooth light fixed pulley at the top of the shaft. To one end of the chain is\par
attached a bucket A of negligible mass and to the other a bucket B of mass m. The system is\par
used to raise ore from the mine as follows. When bucket A is at the top it is filled with mass\par
\[2m of water and bucket B is filled with mass λm of ore, where 0 < λ < 1. The buckets are\]
then released, so that bucket A descends and bucket B ascends. When bucket B reaches the\par
top both buckets are emptied and released, so that bucket B descends and bucket A ascends.\par
The time to fill and empty the buckets is negligible. Find the time taken from the moment\par
bucket A is released at the top until the first time it reaches the top again.\par
This process goes on for a very long time. Show that, if the greatest amount of ore is to be\par
\[raised in that time, then λ must satisfy the condition f (λ) = 0 where\]
\[λ(1 λ)1/2\]
\[f(λ) = - .\]
\[(1 λ)1/2 + (3 + λ)1/2\]
-\par
\end{flushleft}
\subsection*{Page 5}
\begin{flushleft}
\bigskip\textbf{Section C: Probability and Statistics}\par
\bigskip\textbf{Question 12.} Suppose that a solution (X, Y, Z) of the equation\par
\[X + Y + Z = 20,\]
with X, Y and Z non-negative integers, is chosen at random (each such solution being equally\par
likely). Are X and Y independent? Justify your answer.\par
Show that the probability that X is divisible by 5 is 5/21. What is the probability that XY Z\par
is divisible by 5?\par
\bigskip\textbf{Question 13.} I have a bag initially containing r red fruit pastilles (my favourites) and b fruit pastilles of\par
other colours. From time to time I shake the bag thoroughly and remove a pastille at random.\par
(It may be assumed that all pastilles have an equal chance of being selected.) If the pastille\par
is red I eat it but otherwise I replace it in the bag. After n such drawings, I find that I have\par
only eaten one pastille. Show that the probability that I ate it on my last drawing is\par
(r + b\par
1)n-1\par
- .\par
(r + b)n (r + b 1)n\par
- -\par
\bigskip\textbf{Question 14.} To celebrate the opening of the financial year the finance minister of Genland flings a Slihing,\par
a circular coin of radius a cm, where 0 < a < 1, onto a large board divided into squares by\par
two sets of parallel lines 2 cm apart. If the coin does not cross any line, or if the coin covers an\par
intersection, the tax on yaks remains unchanged. Otherwise the tax is doubled. Show that,\par
in order to raise most tax, the value of a should be\par
\[\pi -1\]
1 + .\par
³ ´\par
\[\pi -1\]
\[If, indeed, a = 1 + and the tax on yaks is 1 Slihing per yak this year, show that its\]
³ ´\par
expected value after n years will have passed is\par
n\par
\[8 + \pi\]
.\par
\[µ 4 + \pi ¶\]
\end{flushleft}
\end{document}
